\sect{Task Planning in Robotics}{task}

Task planning forms the reasoning core of autonomous environments, determining sequences of high-level actions that enable a robot to achieve defined objectives under specific constraints. It operates above motion control, focusing on the logical order of operations rather than trajectory computation, and is fundamental for coordinating complex automation systems where multiple agents, sensors, and instruments must act coherently \cite{SicilianoKhatib2008}. Early work adopted AI formalisms such as STRIPS and state-transition systems, in which the world is represented symbolically and actions are described as triplets $\langle name(a), precond(a), effects(a) \rangle$ over a propositional language $\mathcal{L}$, applied under the closed-world assumption \cite{cambon2009hybrid}. While these symbolic approaches provided efficiency and structure, they failed to incorporate geometric and physical constraints, limiting their applicability in real robotic domains where spatial feasibility and uncertainty dominate \cite{cambon2009hybrid}.

To overcome this abstraction gap, semantic reasoning and hybrid task-motion frameworks were introduced. Semantic representations such as knowledge graphs and semantic maps integrate conceptual and spatial knowledge, allowing automated systems to reason about affordances, containment, and relationships among objects \cite{semanticmaps}. Hybrid Task and Motion Planning further merges symbolic reasoning with geometric validation by representing system states as pairs $(s_{sym}, s_{geo})$ and ensuring that both logical and geometric preconditions are satisfied within the configuration space $C_S^{free}$ \cite{cambon2009hybrid}. Probabilistic motion planners like PRM and RRT are used to guarantee reachability and collision-free feasibility. This integration enables autonomous systems to plan compound operations such as pick, place, and measurement actions while preserving consistency between task intent and physical execution.

Despite their sophistication, these frameworks demand extensive manual modeling of objects, grasps, and environment states, which constrains scalability and adaptability in dynamic laboratory or production environments \cite{cambon2009hybrid, semanticmaps}. Consequently, current research focuses on data-driven and AI-based planning paradigms that can infer action structures, interpret under-specified goals, and generalize across diverse setups. Instead of manually encoding every symbolic condition, the emerging approach defines the complete range of robotic actions and movement primitives, allowing an intelligent LLM based agent to autonomously arrange and execute the planning process.